% !TeX root = closed-systems-from-comparison-completeness.tex
% ============================================================
\chapter{Complex Hilbert Structure from Quadratic Transport}
\label{ch:hilbert}
% ============================================================
\section{Introduction}
Under the standing principle of closed-world admissibility
(\cref{sp:closed-world}), this chapter develops the Hilbert-reconstruction
consequences of the transport-visibility clause \cref{sp:transport} on
the stabilized quadratic carrier.
From the quadratic transport carrier, it derives Hilbert structure and
supplies the phase and scalar machinery used in
\cref{ch:einstein,ch:em}.
The comparison target is the usual complex Hilbert-space formalism of
quantum mechanics \cite{VonNeumann1955}; the chapter reconstructs that
structure from the quadratic transport layer rather than assuming it.
\Cref{thm:interface-stabilized,thm:interface-seal} identify the
stabilized quadratic transport
carrier
\[
	\mathcal K \;\simeq\; F^{2}/F^{3}
\]
as the first intrinsic layer at which comparison transport fails to
close strictly.
This carrier is not an auxiliary linearization, nor an externally
imposed state space.
It is the first quotient in the intrinsic transport filtration on which
nontrivial comparison defect survives.
Equivalently, it is the first level of the closed stack at which
transport retains information beyond first-order flatness.
Two structures have already been determined on this quadratic layer.
The first is scalar.
The probability law for mutually exclusive first-visible alternatives is
tracked through their classes in \(\mathcal K\), while, under the
inherited second-jet faithfulness condition, it is determined by the
unique positive quadratic scalar channel on the realized degree--\(2\)
channel attached to the stabilized quadratic carrier.  Thus the
first-visible defect sector is not merely a collection of defect
classes.  Through that realized channel it already carries a
distinguished positive quadratic law.
The second is orientational.
Quadratic defects arise from oriented comparison triangles.
Reversal of triangle orientation reverses the corresponding quadratic
defect.
Hence the same carrier \(\mathcal K\) carries a canonical reversal
symmetry induced by triangle orientation reversal.
The purpose of the present chapter is to show that these two facts
already force complex Hilbert structure.
The subtle point is exact.
An involution by itself does not produce a complex structure:
an operator squaring to the identity is not an operator squaring to
minus the identity.
The missing step is that the quadratic transport carrier is not to be
treated as an unoriented defect space.
Its primitive generators come in two ordered orientation sectors,
related by the unique reversal involution.
Once that ordered orientation double is made explicit, the scalar form
and the reversal symmetry together force a quarter-turn operator,
unique up to sign, and hence a complex structure.
The residual sign ambiguity is structural rather than arbitrary.
It is the necessary symmetry
\[
	J\mapsto -J,
\]
which corresponds to complex conjugation of the same complex Hilbert
structure and not to a second geometric choice.
Once \(J\) is determined, the Hermitian inner product is in turn determined by
the quadratic scalar channel and polarization.
At the end no further datum is chosen.
The resulting classification theorem is the following.
\begin{theorem}[Hilbert structure of the quadratic defect sector]
	\label{thm:hilbert-classification}
	Under the standing principle \cref{sp:closed-world}, let
	\[
		\mathcal K \;\simeq\; F^{2}/F^{3}
	\]
	be the stabilized quadratic transport carrier of the closed comparison
	stack.
	Then the following structures are determined intrinsically:
	\begin{enumerate}[label=\textup{(\roman*)}]
		\item
		      under the inherited second-jet faithfulness condition, a positive
		      definite quadratic scalar form
		      \[
			      Q:\mathcal K_{\mathbb R}\to\mathbb R_{\ge 0},
			      \qquad
			      \mathcal K_{\mathbb R}:=\mathcal K\otimes_{\mathbb Z}\mathbb R,
		      \]
		      extending the canonical degree--\(2\) scalar channel on the
		      realized degree--\(2\) channel attached to the stabilized
		      quadratic carrier;
		\item
		      a unique filtration-compatible orientation involution
		      \[
			      \sigma:\mathcal K_{\mathbb R}\to\mathcal K_{\mathbb R}
		      \]
		      induced by reversal of primitive triangle orientation;
		\item
		      a canonical oriented defect double
		      \[
			      H_{\mathbb R}:=\mathcal K_{\mathbb R}\otimes_{\mathbb R}\mathcal O,
		      \]
		      where \(\mathcal O\) is the real orientation module generated by the
		      two primitive triangle orientations;
		\item
		      an orthogonal complex structure
		      \[
			      J:H_{\mathbb R}\to H_{\mathbb R},
			      \qquad
			      J^{2}=-\mathrm{id},
		      \]
		      unique up to sign among orthogonal endomorphisms compatible with total
		      orientation reversal and with the quadratic transport filtration;
		\item
		      a unique Hermitian inner product on \(H_{\mathbb R}\) whose
		      norm-square extends the quadratic scalar channel.
	\end{enumerate}
	Consequently the completion of \(H_{\mathbb R}\) in the induced norm is
	a complex Hilbert space canonically associated to the quadratic defect
	sector, and, under that same inherited second-jet interface
	hypothesis, the probability law on first-visible alternatives becomes
	the norm-square law through the realized degree--\(2\) channel
	attached to the stabilized quadratic carrier.
\end{theorem}
The proof occupies the remainder of the chapter.
The argument is not a free construction.
Each stage removes a genuine ambiguity.
The quadratic scalar channel determines the real inner product.
The orientation reversal is determined uniquely by the primitive transport
triangles.
The orientation module carries a unique quarter-turn compatible with
reversal, up to overall sign.
The full oriented defect double then carries a unique orthogonal complex
structure anti-commuting with total orientation reversal, again up to
that same harmless sign.
Finally the Hermitian form is determined uniquely by polarization.
At the end nothing remains to be chosen except the sign
\(J\mapsto -J\), and that sign does not represent additional geometry.
% ============================================================
\section{The quadratic defect sector and its scalar form}
\label{sec:hilbert-quadratic-sector}
% ============================================================
Let
\[
	\mathcal K \;\simeq\; F^{2}/F^{3}
\]
denote the stabilized quadratic transport carrier.
\Cref{thm:interface-main,thm:interface-stabilized} show that
\(\mathcal K\) is generated by the quadratic
wedge classes
\[
	[e_i\wedge e_j],
\]
arising from degree--\(2\) commutator residues of oriented comparison
transport.
We first pass from the intrinsic abelian carrier to its real span.
\begin{definition}[Real quadratic defect space]
	Define
	\[
		\mathcal K_{\mathbb R}
		:=
		\mathcal K\otimes_{\mathbb Z}\mathbb R.
	\]
\end{definition}
The scalar-channel theory of
\cref{lem:scalar-reduction,lem:K-scalar-unique,thm:born,thm:quadratic-extension}
yields a canonical
positive quadratic law on this real vector space.
To use it here, one must formulate the extension carefully.
The point is not merely that a scalar law exists on the integral
carrier, but that the previously established quadratic-extension
principle applies to the stabilized degree--\(2\) transport sector.
\begin{theorem}[Quadratic scalar channel]
	\label{thm:quadratic-scalar}
	Under the inherited second-jet faithfulness condition, there exists a
	unique positive quadratic form
	\[
		Q:\mathcal K_{\mathbb R}\to\mathbb R_{\ge 0}
	\]
	extending the canonical degree--\(2\) scalar channel on the
	realized degree--\(2\) channel attached to the stabilized quadratic
	carrier, whose restriction to first-visible alternatives tracked
	through the integral carrier \(\mathcal K\) is the admissible
	branch-weight law.
\end{theorem}
\begin{proof}
	The prior results
	\cref{thm:interface-stabilized,lem:scalar-reduction,lem:K-scalar-unique,thm:quadratic-extension}
	establish three facts.
	First, first-visible alternatives are tracked through the stabilized
	quadratic carrier
	\[
		\mathcal K\simeq F^{2}/F^{3}.
	\]
	Second, under the inherited second-jet faithfulness condition, the
	admissible branch-weight law on those first-visible alternatives is
	the canonical degree--\(2\) scalar channel on the realized
	degree--\(2\) channel attached to that stabilized quadratic carrier,
	unique up to overall normalization among degree--\(2\) scalarizations.
	Third, the quadratic-extension theorem applies to that realized-channel
	scalar law and extends it uniquely to the real linearization
	\(\mathcal K_{\mathbb R}=\mathcal K\otimes_{\mathbb Z}\mathbb R\).
	Applying that extension theorem in the present inherited-interface
	regime yields a unique quadratic form
	\[
		Q:\mathcal K_{\mathbb R}\to\mathbb R_{\ge 0}
	\]
	whose restriction to first-visible alternatives tracked through
	\(\mathcal K\) is the admissible branch-weight law.  Because that law
	is nonnegative on first-visible alternatives, the extended quadratic
	form is nonnegative as well.  This is the required map \(Q\).
\end{proof}
\begin{remark}[What is used here]
	\Cref{thm:quadratic-scalar} does not introduce a new scalar law.
	It imports, under the inherited second-jet faithfulness condition, the
	canonical degree--\(2\) scalar channel on the realized degree--\(2\)
	channel attached to the stabilized quadratic carrier and extends
	that law to the real span \(\mathcal K_{\mathbb R}\).
	Thus the scalar geometry entering the Hilbert construction is already
	present in the transport analysis of
	\cref{thm:interface-stabilized,thm:born,thm:quadratic-extension}.
\end{remark}
The decisive point is that positivity is strict, not merely weak.
\begin{lemma}[Strict positivity of the quadratic scalar channel]
	\label{lem:hilbert-strict-positivity}
	For every nonzero vector
	\[
		x\in \mathcal K_{\mathbb R},
	\]
	one has
	\[
		Q(x)>0.
	\]
	Equivalently,
	\[
		Q(x)=0 \Longrightarrow x=0.
	\]
\end{lemma}
\begin{proof}
	We first prove the claim on the integral carrier \(\mathcal K\).
	Let
	\[
		0\neq k\in \mathcal K.
	\]
	Since \(\mathcal K\simeq F^{2}/F^{3}\) is the stabilized first visible
	transport-defect carrier, a nonzero class \(k\) is, by definition, a
	nontrivial first-visible defect class.
	The scalar-channel theorem identifies the admissible branch-weight law
	as the unique positive scalar detector of such first-visible classes.
	Hence
	\[
		Q(k)>0
		\qquad
		\text{for every }0\neq k\in\mathcal K.
	\]
	Now let
	\[
		0\neq x\in \mathcal K_{\mathbb R}.
	\]
	Choose an integral basis
	\[
		k_1,\dots,k_m
	\]
	of the free part of \(\mathcal K\), so that
	\[
		x=\sum_{a=1}^{m} \lambda_a (k_a\otimes 1)
	\]
	with real coefficients \(\lambda_a\), not all zero.
	By density of rational coefficients and homogeneity of degree \(2\), it
	suffices to prove strict positivity on nonzero rational combinations.
	Thus choose \(N>0\) such that
	\[
		Nx\in \mathcal K
	\]
	and \(Nx\neq 0\).
	Then, by the integral case already proved,
	\[
		Q(Nx)>0.
	\]
	Since \(Q\) is quadratic,
	\[
		Q(Nx)=N^{2}Q(x).
	\]
	Therefore
	\[
		Q(x)=N^{-2}Q(Nx)>0.
	\]
	This proves strict positivity on all nonzero vectors of
	\(\mathcal K_{\mathbb R}\).
\end{proof}
The quadratic form determines a symmetric bilinear form by polarization.
\begin{definition}[Real polarization]
	Define
	\[
		g_{\mathcal K}(x,y)
		:=
		\frac12\bigl(Q(x+y)-Q(x)-Q(y)\bigr)
		\qquad
		(x,y\in\mathcal K_{\mathbb R}).
	\]
\end{definition}
\begin{proposition}[Real inner product on the quadratic carrier]
	\label{prop:hilbert-real-inner}
	The form \(g_{\mathcal K}\) is a positive definite symmetric bilinear
	form on \(\mathcal K_{\mathbb R}\).
\end{proposition}
\begin{proof}
	Because \(Q\) is a quadratic form on the real vector space
	\(\mathcal K_{\mathbb R}\), there exists a unique symmetric bilinear
	form whose diagonal is \(Q\), and the displayed formula is exactly its
	polarization.
	Thus \(g_{\mathcal K}\) is symmetric and bilinear.
	Moreover, for every \(x\in\mathcal K_{\mathbb R}\),
	\[
		g_{\mathcal K}(x,x)=Q(x)\ge 0.
	\]
	If
	\[
		g_{\mathcal K}(x,x)=0,
	\]
	then
	\[
		Q(x)=0,
	\]
	and \cref{lem:hilbert-strict-positivity} implies
	\[
		x=0.
	\]
	Therefore \(g_{\mathcal K}\) is positive definite.
\end{proof}
\begin{remark}[No complex structure yet]
	At this stage the quadratic carrier has acquired a canonical real
	Euclidean geometry and nothing more.
	The scalar form determines \(g_{\mathcal K}\), but a real inner product
	alone does not force complex structure.
	The next task is to isolate the missing orientational input.
\end{remark}
% ============================================================
\section{Orientation reversal on the quadratic carrier}
\label{sec:hilbert-orientation}
% ============================================================
The next step is to isolate the orientational structure determined by
comparison triangles.
\begin{lemma}[Orientation involution]
	\label{lem:orientation-involution}
	Primitive triangle orientation reversal induces a linear involution
	\[
		\sigma:\mathcal K_{\mathbb R}\to\mathcal K_{\mathbb R}
	\]
	satisfying
	\[
		\sigma^{2}=\mathrm{id},
		\qquad
		\sigma([e_i\wedge e_j])=-[e_i\wedge e_j].
	\]
\end{lemma}
\begin{proof}
	Let \((i,j,\ell)\) be an oriented primitive comparison triangle.
	Reversal of orientation sends \((i,j,\ell)\) to \((\ell,j,i)\).
	At the level of transport words, this reverses the traversal order of
	the corresponding elementary comparison steps.
	Passing to the quadratic quotient
	\[
		F^{2}/F^{3},
	\]
	reversal of order changes the sign of the associated wedge class.
	Thus
	\[
		[e_i\wedge e_j]\longmapsto -[e_i\wedge e_j].
	\]
	Because primitive triangle reversal is itself an involution, the induced
	map on \(\mathcal K\), and hence on \(\mathcal K_{\mathbb R}\), squares
	to the identity.
	Linearity is immediate from passage to the real span.
\end{proof}
\begin{lemma}[Orthogonality of the orientation involution]
	\label{lem:orientation-orthogonal}
	The involution \(\sigma\) preserves the quadratic scalar form and is
	orthogonal with respect to \(g_{\mathcal K}\):
	\[
		Q(\sigma x)=Q(x),
		\qquad
		g_{\mathcal K}(\sigma x,\sigma y)=g_{\mathcal K}(x,y).
	\]
\end{lemma}
\begin{proof}
	Orientation reversal changes a primitive quadratic defect to its
	negative.
	Since the scalar channel is quadratic, it is invariant under sign:
	\[
		Q(-x)=Q(x)
		\qquad
		(x\in\mathcal K_{\mathbb R}).
	\]
	By \cref{lem:orientation-involution}, \(\sigma\) acts on primitive
	generators by sign reversal, hence
	\[
		Q(\sigma x)=Q(x)
	\]
	on the spanning set of primitive generators, and therefore on all of
	\(\mathcal K_{\mathbb R}\) by quadraticity and linear extension.
	Orthogonality of \(\sigma\) with respect to \(g_{\mathcal K}\) now
	follows by polarization:
	\[
		g_{\mathcal K}(\sigma x,\sigma y)
		=
		\frac12\bigl(Q(\sigma x+\sigma y)-Q(\sigma x)-Q(\sigma y)\bigr)
		=
		\frac12\bigl(Q(\sigma(x+y))-Q(x)-Q(y)\bigr)
		=
		g_{\mathcal K}(x,y).
	\]
\end{proof}
At this point one has an involution, but not yet uniqueness.
That distinction matters.
The later complex structure must not depend on a discretionary choice of
reversal operator.
One must show that the transport filtration itself determines the
orientation involution and excludes all competitors.
\begin{lemma}[Uniqueness of the orientation involution]
	\label{lem:orientation-unique}
	Let
	\[
		\tau:\mathcal K_{\mathbb R}\to\mathcal K_{\mathbb R}
	\]
	be a linear automorphism satisfying the following properties:
	\begin{enumerate}[label=\textup{(\roman*)}]
		\item
		      \(\tau\) is induced by an automorphism of the primitive triangle
		      transport system preserving the quadratic filtration and passage to
		      \(F^{2}/F^{3}\);
		\item
		      \(\tau\) reverses primitive triangle orientation;
		\item
		      \(\tau^{2}=\mathrm{id}\).
	\end{enumerate}
	Then
	\[
		\tau=\sigma.
	\]
\end{lemma}
\begin{proof}
	The quadratic carrier \(\mathcal K_{\mathbb R}\) is spanned by the wedge
	classes
	\[
		[e_i\wedge e_j].
	\]
	Because \(\tau\) is filtration-compatible, its action on the entire
	carrier is determined by its action on this generating family.
	Now \(\tau\) reverses primitive triangle orientation.
	Therefore on each primitive wedge generator it acts exactly as
	orientation reversal acts:
	\[
		\tau([e_i\wedge e_j])=-[e_i\wedge e_j].
	\]
	By \cref{lem:orientation-involution}, the same formula holds for
	\(\sigma\).
	Hence \(\tau\) and \(\sigma\) agree on a spanning set of
	\(\mathcal K_{\mathbb R}\), and therefore agree everywhere.
\end{proof}
\begin{remark}[Rigidity of reversal]
	The involution \(\sigma\) is therefore not merely available.
	It is the unique filtration-compatible realization of triangle
	orientation reversal on the quadratic carrier.
	This rigidity is what prevents the later complex structure from hiding a
	new arbitrary choice.
\end{remark}
\subsection{Uniqueness of orientation splitting on the quadratic carrier}
\label{sec:orientation-unique}
The previous lemmas can now be assembled into the object-locus
uniqueness statement used later in \cref{ch:matter}.
\begin{theorem}[Uniqueness of orientation splitting]
\label{thm:orientation-splitting-uniqueness}
Under the standing principle \cref{sp:closed-world}, there exists a
nontrivial orientation involution
\[
	R:\mathcal K_{\mathbb R}\to \mathcal K_{\mathbb R},
	\qquad
	R^2=\id,
\]
reversing primitive oriented quadratic observables.
Moreover, up to conjugacy by admissible filtration-compatible
automorphisms of \(\mathcal K_{\mathbb R}\), this involution is unique.
Its nontrivial action determines a canonical minimal irreducible real
block on which there is an orthogonal quarter-turn operator
\(J\) satisfying
\[
	J^2=-\id,
	\qquad
	JR=-RJ.
\]
\end{theorem}
\begin{proof}
	\textbf{Step 1 (transport visibility: \cref{sp:transport}).}
	Existence of a nontrivial orientation involution is
	\cref{lem:orientation-involution}.
	Orthogonality and filtration compatibility are
	\cref{lem:orientation-orthogonal}.

	\textbf{Step 2 (intrinsic and quotient admissibility:
	\cref{sp:intrinsic,sp:quotient}).}
	Uniqueness among admissible orientation-reversing involutions is
	\cref{lem:orientation-unique}.
	Thus the orientation reversal channel is fixed up to admissible
	filtration-compatible equivalence.

	\textbf{Step 3 (intrinsic minimal block extraction:
	\cref{sp:intrinsic}).}
	Let \(V\subseteq\mathcal K_{\mathbb R}\) be a minimal nontrivial
	\(R\)-stable real subspace.
	Pick \(0\neq v\in V\).
	Since \(R\) is nontrivial, \(Rv\) is not proportional to \(v\), and
	minimality gives
	\[
		V=\mathrm{span}_{\mathbb R}\{v,Rv\}.
	\]
	Define
	\[
		Jv:=Rv,
		\qquad
		J(Rv):=-v,
	\]
	and extend linearly.
	Then
	\[
		J^2=-\id,
		\qquad
		JR=-RJ.
	\]
	Because \(R\) is orthogonal and \(V\) is generated by an
	\(R\)-orbit pair, \(J\) is orthogonal on \(V\).
	Conjugacy uniqueness follows from uniqueness of the involution and
	filtration-compatible transport equivalence.
\end{proof}
\begin{corollary}[Primitive doublet]
\label{cor:orientation-primitive-doublet}
The minimal irreducible real representation carrying the pair
\((R,J)\) of \cref{thm:orientation-splitting-uniqueness} is
two-dimensional.
Hence orientation splitting determines a primitive multiplicity generator
of dimension \(2\).
\end{corollary}
\begin{proof}
On the minimal nontrivial block, \(R\) has eigenvalues \(\pm1\) and
\(J\) exchanges the eigendirections because \(JR=-RJ\).
Thus at least two real dimensions are required.
Minimality of the block excludes larger dimension.
\end{proof}
% ============================================================
\section{The orientation module}
\label{sec:hilbert-orientation-module}
% ============================================================
An involution on \(\mathcal K_{\mathbb R}\) does not by itself produce a
complex structure.
The relevant object is not the bare carrier
\(\mathcal K_{\mathbb R}\), but the ordered double of its two primitive
orientation sectors.
We now make that double explicit.
\begin{definition}[Orientation module]
	\label{def:orientation-module}
	Let \(\mathcal O\) be the two-dimensional real vector space with basis
	\[
		\varepsilon_{+},\ \varepsilon_{-},
	\]
	representing the two primitive triangle orientations.
	Equip \(\mathcal O\) with the Euclidean inner product
	\(h_{\mathcal O}\) for which this basis is orthonormal, and define the
	reversal involution
	\[
		r:\mathcal O\to\mathcal O
	\]
	by
	\[
		r(\varepsilon_{+})=\varepsilon_{-},
		\qquad
		r(\varepsilon_{-})=\varepsilon_{+}.
	\]
\end{definition}
The key point is that \(\mathcal O\) carries a quarter-turn compatible
with reversal, and that this quarter-turn is unique up to sign.
\begin{proposition}[Quarter-turn on the orientation module]
	\label{prop:orientation-quarter-turn}
	There exists an orthogonal endomorphism
	\[
		j_{\mathcal O}:\mathcal O\to\mathcal O
	\]
	such that
	\[
		j_{\mathcal O}^{2}=-\mathrm{id},
		\qquad
		j_{\mathcal O}r=-rj_{\mathcal O}.
	\]
	It is unique up to overall sign.
	Explicitly,
	\[
		j_{\mathcal O}(\varepsilon_{+})=\varepsilon_{-},
		\qquad
		j_{\mathcal O}(\varepsilon_{-})=-\varepsilon_{+}.
	\]
\end{proposition}
\begin{proof}
	Define \(j_{\mathcal O}\) by the displayed formulas.
	Then
	\[
		j_{\mathcal O}^{2}(\varepsilon_{+})
		=
		j_{\mathcal O}(\varepsilon_{-})
		=
		-\varepsilon_{+},
		\qquad
		j_{\mathcal O}^{2}(\varepsilon_{-})
		=
		j_{\mathcal O}(-\varepsilon_{+})
		=
		-\varepsilon_{-},
	\]
	so
	\[
		j_{\mathcal O}^{2}=-\mathrm{id}.
	\]
	In the ordered orthonormal basis
	\(\{\varepsilon_{+},\varepsilon_{-}\}\), the matrix of \(j_{\mathcal O}\)
	is
	\[
		\begin{pmatrix}
			0 & -1 \\
			1 & 0
		\end{pmatrix},
	\]
	hence \(j_{\mathcal O}\) is orthogonal.
	Next,
	\[
		j_{\mathcal O}r(\varepsilon_{+})
		=
		j_{\mathcal O}(\varepsilon_{-})
		=
		-\varepsilon_{+},
		\qquad
		-rj_{\mathcal O}(\varepsilon_{+})
		=
		-r(\varepsilon_{-})
		=
		-\varepsilon_{+},
	\]
	and similarly on \(\varepsilon_{-}\), so
	\[
		j_{\mathcal O}r=-rj_{\mathcal O}.
	\]
	For uniqueness, let \(j'_{\mathcal O}\) be another orthogonal
	endomorphism satisfying
	\[
		(j'_{\mathcal O})^{2}=-\mathrm{id},
		\qquad
		j'_{\mathcal O}r=-rj'_{\mathcal O}.
	\]
	The operator \(r\) has eigenspaces spanned by
	\[
		\varepsilon_{+}+\varepsilon_{-}
		\qquad\text{and}\qquad
		\varepsilon_{+}-\varepsilon_{-}.
	\]
	The anti-commutation relation determines \(j'_{\mathcal O}\) to exchange
	these two lines.
	Orthogonality and the identity
	\((j'_{\mathcal O})^{2}=-\mathrm{id}\) then leave exactly two
	possibilities:
	\[
		j'_{\mathcal O}=j_{\mathcal O}
		\qquad\text{or}\qquad
		j'_{\mathcal O}=-j_{\mathcal O}.
	\]
\end{proof}
\begin{remark}[Meaning of the sign ambiguity]
	The sign ambiguity in \(j_{\mathcal O}\) is necessary and harmless.
	It is the ordinary ambiguity between a complex structure and its
	conjugate.
	It does not encode a second orientation geometry on the closed stack.
\end{remark}
% ============================================================
\section{The oriented defect double and its universal property}
\label{sec:hilbert-oriented-double}
% ============================================================
We now form the real vector space on which the complex structure will
live.
\begin{definition}[Oriented defect double]
	Define
	\[
		H_{\mathbb R}:=\mathcal K_{\mathbb R}\otimes_{\mathbb R}\mathcal O.
	\]
	Equip \(H_{\mathbb R}\) with the symmetric bilinear form
	\[
		g(x\otimes u,\,y\otimes v)
		:=
		g_{\mathcal K}(x,y)\,h_{\mathcal O}(u,v),
	\]
	extended bilinearly to all of \(H_{\mathbb R}\).
\end{definition}
\begin{proposition}[Product metric]
	The form \(g\) is a positive definite symmetric bilinear form on
	\(H_{\mathbb R}\).
\end{proposition}
\begin{proof}
	Symmetry and bilinearity are immediate from symmetry and bilinearity of
	\(g_{\mathcal K}\) and \(h_{\mathcal O}\).
	It remains to prove positive definiteness.
	Let
	\[
		\xi=\sum_{a=1}^{m} x_a\otimes u_a\in H_{\mathbb R}.
	\]
	Choose the orthonormal basis
	\[
		\varepsilon_{+},\varepsilon_{-}
	\]
	of \(\mathcal O\) from \cref{def:orientation-module}.
	Then \(\xi\) may be written uniquely in the form
	\[
		\xi=x_{+}\otimes \varepsilon_{+}+x_{-}\otimes \varepsilon_{-}
	\]
	for some \(x_{+},x_{-}\in\mathcal K_{\mathbb R}\).
	Since the basis is orthonormal, one has
	\[
		g(\xi,\xi)
		=
		g_{\mathcal K}(x_{+},x_{+})+g_{\mathcal K}(x_{-},x_{-}).
	\]
	By \cref{prop:hilbert-real-inner}, the form \(g_{\mathcal K}\) is
	positive definite.
	Hence
	\[
		g(\xi,\xi)\ge 0,
	\]
	with equality if and only if
	\[
		x_{+}=0
		\qquad\text{and}\qquad
		x_{-}=0.
	\]
	Thus \(\xi=0\).
	Therefore \(g\) is positive definite.
\end{proof}
The two sources of involution on \(H_{\mathbb R}\) are the quadratic
orientation involution \(\sigma\) on \(\mathcal K_{\mathbb R}\) and the
sector-reversal involution \(r\) on \(\mathcal O\).
Together they determine the global orientation-reversal symmetry.
\begin{definition}[Total orientation reversal]
	Define
	\[
		\Sigma:=\sigma\otimes r
		:
		H_{\mathbb R}\to H_{\mathbb R}.
	\]
\end{definition}
\begin{proposition}[Total reversal is orthogonal]
	The operator \(\Sigma\) is an orthogonal involution:
	\[
		\Sigma^{2}=\mathrm{id},
		\qquad
		g(\Sigma\xi,\Sigma\eta)=g(\xi,\eta)
		\qquad
		(\xi,\eta\in H_{\mathbb R}).
	\]
\end{proposition}
\begin{proof}
	Since
	\[
		\sigma^{2}=\mathrm{id}_{\mathcal K_{\mathbb R}}
		\qquad\text{and}\qquad
		r^{2}=\mathrm{id}_{\mathcal O},
	\]
	one has
	\[
		\Sigma^{2}
		=
		(\sigma\otimes r)(\sigma\otimes r)
		=
		\sigma^{2}\otimes r^{2}
		=
		\mathrm{id}_{H_{\mathbb R}}.
	\]
	To prove orthogonality, it suffices to check the formula on pure
	tensors.
	For \(x,y\in\mathcal K_{\mathbb R}\) and \(u,v\in\mathcal O\),
	\[
		g(\Sigma(x\otimes u),\Sigma(y\otimes v))
		=
		g(\sigma x\otimes ru,\sigma y\otimes rv)
		=
		g_{\mathcal K}(\sigma x,\sigma y)\,h_{\mathcal O}(ru,rv).
	\]
	By \cref{lem:orientation-orthogonal},
	\[
		g_{\mathcal K}(\sigma x,\sigma y)=g_{\mathcal K}(x,y),
	\]
	and by construction of \(r\),
	\[
		h_{\mathcal O}(ru,rv)=h_{\mathcal O}(u,v).
	\]
	Therefore
	\[
		g(\Sigma(x\otimes u),\Sigma(y\otimes v))
		=
		g_{\mathcal K}(x,y)\,h_{\mathcal O}(u,v)
		=
		g(x\otimes u,y\otimes v).
	\]
	By bilinearity this holds for all \(\xi,\eta\in H_{\mathbb R}\).
\end{proof}
The appropriate quarter-turn on \(H_{\mathbb R}\) is now determined by the
orientation-module quarter-turn.
The only remaining issue is to exclude hidden carrier automorphisms.
\begin{lemma}[Rigidity of filtration-compatible carrier automorphisms]
	\label{lem:hilbert-carrier-rigidity}
	Let
	\[
		A:\mathcal K_{\mathbb R}\to\mathcal K_{\mathbb R}
	\]
	be an orthogonal automorphism satisfying the following properties:
	\begin{enumerate}[label=\textup{(\roman*)}]
		\item
		      \(A\) is induced functorially from the quadratic transport filtration;
		\item
		      \(A\) commutes with the orientation involution:
		      \[
			      A\sigma=\sigma A;
		      \]
		\item
		      \(A\) preserves the primitive wedge generator system up to filtration
		      equivalence.
	\end{enumerate}
	Then
	\[
		A=\pm\mathrm{id}_{\mathcal K_{\mathbb R}}.
	\]
	If in addition \(A\) preserves primitive orientation sector by sector,
	then
	\[
		A=\mathrm{id}_{\mathcal K_{\mathbb R}}.
	\]
\end{lemma}
\begin{proof}
	The carrier \(\mathcal K_{\mathbb R}\) is generated by the primitive
	quadratic wedge classes
	\[
		[e_i\wedge e_j].
	\]
	By filtration-functoriality and preservation of the primitive generator
	system up to filtration equivalence, the action of \(A\) is determined
	by its action on these generators.
	Since \(A\) commutes with \(\sigma\) and \(\sigma\) acts on each
	primitive generator by sign reversal,
	\[
		\sigma([e_i\wedge e_j])=-[e_i\wedge e_j],
	\]
	the image under \(A\) of a primitive generator must again lie in the
	same one-dimensional orientation line generated by
	\([e_i\wedge e_j]\).
	Indeed, if \(A\) mixed distinct primitive lines, then commuting with the
	unique orientation involution would produce an additional
	filtration-compatible intrinsic symmetry on the primitive quadratic
	transport carrier beyond the determined reversal symmetry, contradicting the
	rigidity of primitive transport classification in
	\cref{thm:hilbert-classification,lem:orientation-unique}.
	Therefore, for each primitive generator,
	\[
		A([e_i\wedge e_j])=\lambda_{ij}[e_i\wedge e_j]
	\]
	for some real scalar \(\lambda_{ij}\).
	Because \(A\) is orthogonal with respect to the positive definite form
	\(g_{\mathcal K}\), each \(\lambda_{ij}\) satisfies
	\[
		\lambda_{ij}^{2}=1,
	\]
	hence
	\[
		\lambda_{ij}=\pm 1.
	\]
	It remains to show that the sign is global rather than generatorwise.
	If different primitive generators carried different signs, then \(A\)
	would separate distinct primitive transport directions by an additional
	filtration-compatible orthogonal datum not already encoded by the
	quadratic transport system.
	But the quadratic rigidity results
	(\cref{thm:hilbert-classification,lem:orientation-unique}) identify the primitive
	degree--\(2\) carrier only through the scalar form and the unique
	orientation involution.
	Accordingly no further generatorwise sign freedom exists.
	Hence all \(\lambda_{ij}\) coincide, and therefore
	\[
		A=\pm \mathrm{id}_{\mathcal K_{\mathbb R}}.
	\]
	If \(A\) preserves primitive orientation sector by sector, then the
	global negative sign is excluded, and one obtains
	\[
		A=\mathrm{id}_{\mathcal K_{\mathbb R}}.
	\]
\end{proof}
\begin{theorem}[Universal complex structure on the oriented defect double]
	\label{thm:universal-complex-structure}
	There exists an orthogonal endomorphism
	\[
		J:H_{\mathbb R}\to H_{\mathbb R}
	\]
	such that
	\[
		J^{2}=-\mathrm{id},
		\qquad
		J\Sigma=-\Sigma J.
	\]
	It is unique up to overall sign among orthogonal endomorphisms with
	these properties and compatible with the quadratic transport
	filtration.
	Explicitly,
	\[
		J:=\mathrm{id}_{\mathcal K_{\mathbb R}}\otimes j_{\mathcal O}.
	\]
\end{theorem}
\begin{proof}
	Define
	\[
		J:=\mathrm{id}_{\mathcal K_{\mathbb R}}\otimes j_{\mathcal O}.
	\]
	Then
	\[
		J^{2}
		=
		\mathrm{id}_{\mathcal K_{\mathbb R}}\otimes j_{\mathcal O}^{2}
		=
		-\mathrm{id}_{H_{\mathbb R}},
	\]
	because \(j_{\mathcal O}^{2}=-\mathrm{id}_{\mathcal O}\) by
	\cref{prop:orientation-quarter-turn}.
	Orthogonality follows immediately from orthogonality of
	\(j_{\mathcal O}\) and orthogonality of the identity on
	\(\mathcal K_{\mathbb R}\).
	Next,
	\[
		J\Sigma
		=
		(\mathrm{id}\otimes j_{\mathcal O})(\sigma\otimes r)
		=
		\sigma\otimes (j_{\mathcal O}r)
		=
		-\sigma\otimes (rj_{\mathcal O})
		=
		-\Sigma J,
	\]
	because \(j_{\mathcal O}r=-rj_{\mathcal O}\).
	This proves existence.
	We now prove uniqueness up to sign.
	Let
	\[
		J':H_{\mathbb R}\to H_{\mathbb R}
	\]
	be an orthogonal endomorphism satisfying
	\[
		(J')^{2}=-\mathrm{id},
		\qquad
		J'\Sigma=-\Sigma J',
	\]
	and assume \(J'\) is compatible with the quadratic transport filtration.
	Because \(J'\) is filtration-compatible, its action on the
	\(\mathcal K_{\mathbb R}\)-factor is induced from a
	filtration-compatible orthogonal automorphism
	\[
		A:\mathcal K_{\mathbb R}\to\mathcal K_{\mathbb R}
	\]
	commuting with the unique orientation involution \(\sigma\).
	By \cref{lem:hilbert-carrier-rigidity}, such an \(A\) equals
	\[
		A=\pm \mathrm{id}_{\mathcal K_{\mathbb R}}.
	\]
	Absorbing this harmless global sign into the orientation factor, one may
	therefore write
	\[
		J'= \mathrm{id}_{\mathcal K_{\mathbb R}}\otimes j'_{\mathcal O}
	\]
	for some orthogonal endomorphism
	\[
		j'_{\mathcal O}:\mathcal O\to\mathcal O.
	\]
	The identities for \(J'\) now become
	\[
		(j'_{\mathcal O})^{2}=-\mathrm{id},
		\qquad
		j'_{\mathcal O}r=-rj'_{\mathcal O}.
	\]
	By \cref{prop:orientation-quarter-turn}, the only orthogonal
	endomorphisms of \(\mathcal O\) satisfying these identities are
	\[
		j'_{\mathcal O}=j_{\mathcal O}
		\qquad\text{or}\qquad
		j'_{\mathcal O}=-j_{\mathcal O}.
	\]
	Therefore
	\[
		J'=\pm J.
	\]
	This proves uniqueness up to sign.
\end{proof}
\begin{remark}[Where complex structure first appears]
	\Cref{thm:universal-complex-structure} is the precise point at which
	complex structure
	becomes necessary.
	The scalar geometry on \(\mathcal K_{\mathbb R}\) supplies the real
	metric.
	The ordered orientation double supplies the quarter-turn.
	The compatibility condition with total orientation reversal removes all
	remaining freedom except the necessary sign \(J\mapsto -J\).
\end{remark}
\begin{proposition}[Uniqueness of the orthogonal complex structure]
	\label{prop:hilbert-J-unique}
	Let
	\[
		H_{\mathbb R}=\mathcal K_{\mathbb R}\otimes_{\mathbb R}\mathcal O
	\]
	be the oriented defect double, equipped with the product inner product
	\(g\), and let
	\[
		\Sigma=\sigma\otimes r
	\]
	be the total orientation-reversal involution.
	Suppose
	\[
		J':H_{\mathbb R}\to H_{\mathbb R}
	\]
	is an orthogonal endomorphism satisfying
	\[
		(J')^{2}=-\mathrm{id},
		\qquad
		J'\Sigma=-\Sigma J',
	\]
	and is induced functorially from the quadratic transport carrier.
	Then
	\[
		J'=\pm J,
	\]
	where
	\[
		J=\mathrm{id}_{\mathcal K_{\mathbb R}}\otimes j_{\mathcal O}
	\]
	is the canonical complex structure from
	\cref{thm:universal-complex-structure}.
\end{proposition}
\begin{proof}
	This is the uniqueness statement already proved in
	\cref{thm:universal-complex-structure}, stated separately for later
	reference.
	The filtration-compatible functoriality condition reduces the carrier
	part to the identity up to harmless global sign by
	\cref{lem:hilbert-carrier-rigidity}, and the orientation quarter-turn is
	unique up to sign by \cref{prop:orientation-quarter-turn}.
	Therefore \(J'=\pm J\).
\end{proof}
In particular, once filtration compatibility and anti-commutation with
total orientation reversal are imposed, the complex structure on the
oriented defect double is determined uniquely up to the harmless sign
ambiguity
\[
	J\mapsto -J.
\]
This sign is structural.
It does not encode a second geometry; it is the analogue of passing from
a complex structure to its conjugate.
% ============================================================
\section{Skew form and Hermitian structure}
\label{sec:hilbert-skew-hermitian}
% ============================================================
The canonical complex structure \(J\) determines a corresponding skew
form.
\begin{definition}[Canonical skew form]
	Define
	\[
		\omega(\xi,\eta):=g(\xi,J\eta)
		\qquad
		(\xi,\eta\in H_{\mathbb R}).
	\]
\end{definition}
\begin{proposition}[Properties of the skew form]
	The form \(\omega\) is bilinear, antisymmetric, and nondegenerate.
\end{proposition}
\begin{proof}
	Bilinearity is immediate from bilinearity of \(g\) and linearity of
	\(J\).
	To prove antisymmetry, we first note that \(J\) is skew-adjoint with
	respect to \(g\).
	Indeed, orthogonality of \(J\) and the identity \(J^{2}=-\mathrm{id}\)
	give, for all \(\xi,\eta\in H_{\mathbb R}\),
	\[
		g(J\xi,\eta)
		=
		g(J^{2}\xi,J\eta)
		=
		-g(\xi,J\eta).
	\]
	Hence
	\[
		\omega(\eta,\xi)
		=
		g(\eta,J\xi)
		=
		-g(J\eta,\xi)
		=
		-g(\xi,J\eta)
		=
		-\omega(\xi,\eta).
	\]
	To prove nondegeneracy, suppose \(\omega(\xi,\eta)=0\) for all
	\(\eta\in H_{\mathbb R}\).
	Then
	\[
		g(\xi,J\eta)=0
		\qquad
		\text{for all }\eta.
	\]
	Since \(J\) is invertible, every vector \(\zeta\in H_{\mathbb R}\) may
	be written as \(\zeta=J\eta\) for some \(\eta\).
	Therefore
	\[
		g(\xi,\zeta)=0
		\qquad
		\text{for all }\zeta\in H_{\mathbb R}.
	\]
	Because \(g\) is nondegenerate, this determines \(\xi=0\).
	Thus \(\omega\) is nondegenerate.
\end{proof}
We now isolate the rigidity statement that removes the last remaining
loophole in the passage from the real metric to the Hermitian form.
\begin{lemma}[Polarization rigidity]
	\label{lem:polarization-rigidity}
	Let \(H_{\mathbb R}\) be the real oriented defect sector, let
	\[
		Q_H(\xi):=g(\xi,\xi)
	\]
	be the quadratic form induced by \(g\), and let
	\[
		J:H_{\mathbb R}\to H_{\mathbb R}
	\]
	be the canonical complex structure of
	\cref{thm:universal-complex-structure}.
	Then there exists a unique Hermitian inner product
	\[
		\langle\cdot,\cdot\rangle:H_{\mathbb R}\times H_{\mathbb R}\to\mathbb C
	\]
	whose real part is \(g\), whose imaginary part is compatible with \(J\),
	and whose norm-square satisfies
	\[
		\langle \xi,\xi\rangle = Q_H(\xi)
		\qquad
		(\xi\in H_{\mathbb R}).
	\]
	More explicitly, the unique such Hermitian form is
	\[
		\langle \xi,\eta\rangle
		=
		g(\xi,\eta)+ i\, g(\xi,J\eta).
	\]
\end{lemma}
\begin{proof}
	We first prove uniqueness.
	Let
	\[
		\langle\cdot,\cdot\rangle
	\]
	be any Hermitian form satisfying the stated properties.
	Its real part is determined uniquely by the norm-square \(Q_H\), because
	real polarization gives
	\[
		\Re\langle \xi,\eta\rangle
		=
		\frac12\bigl(Q_H(\xi+\eta)-Q_H(\xi)-Q_H(\eta)\bigr)
		=
		g(\xi,\eta).
	\]
	Thus the real part is determined.
	The compatibility with \(J\) determines the imaginary part as well.
	Indeed, by definition of compatibility, the imaginary part must be the
	skew form determined by \(g\) and \(J\), namely
	\[
		\Im\langle \xi,\eta\rangle
		=
		g(\xi,J\eta).
	\]
	Hence the Hermitian form is uniquely determined and must equal
	\[
		g(\xi,\eta)+i\,g(\xi,J\eta).
	\]
	It remains to verify that this displayed formula defines a Hermitian
	inner product.
	Conjugate symmetry follows from symmetry of \(g\) and skew-adjointness
	of \(J\):
	\[
		\overline{\langle \xi,\eta\rangle}
		=
		g(\xi,\eta)-i\,g(\xi,J\eta)
		=
		g(\eta,\xi)+i\,g(\eta,J\xi)
		=
		\langle \eta,\xi\rangle.
	\]
	Real bilinearity of \(g\) and linearity of \(J\) imply complex
	sesquilinearity in the usual way when \(H_{\mathbb R}\) is regarded as a
	complex vector space via \(J\).
	Finally, for every \(\xi\in H_{\mathbb R}\),
	\[
		\langle \xi,\xi\rangle
		=
		g(\xi,\xi)+i\,g(\xi,J\xi).
	\]
	But antisymmetry of the skew form gives
	\[
		g(\xi,J\xi)=\omega(\xi,\xi)=0.
	\]
	Therefore
	\[
		\langle \xi,\xi\rangle
		=
		g(\xi,\xi)
		=
		Q_H(\xi)\ge 0.
	\]
	If \(\langle \xi,\xi\rangle=0\), then \(Q_H(\xi)=0\), and since \(g\) is
	positive definite this determines \(\xi=0\).
	Thus the form is positive definite.
	Hence the displayed formula defines the unique Hermitian inner product
	with the required properties.
\end{proof}
\begin{remark}[No residual freedom]
	Once the quadratic scalar law and the canonical complex structure are
	fixed, no further freedom remains.
	The Hermitian inner product is not chosen.
	It is rigidly determined by the real metric and the quarter-turn through
	polarization.
\end{remark}
% ============================================================
\section{Hilbert reconstruction}
\label{sec:hilbert-reconstruction}
% ============================================================
We may now formulate the Hilbert reconstruction theorem in its final
form.
\begin{theorem}[Hilbert reconstruction]
	The quadratic scalar channel \(Q\) on the quadratic carrier and the
	canonical complex structure \(J\) on the oriented defect double
	determine a unique Hermitian inner product on \(H_{\mathbb R}\), namely
	the form of \cref{lem:polarization-rigidity}.
	The completion of \(H_{\mathbb R}\) in the associated norm is a complex
	Hilbert space.
\end{theorem}
\begin{proof}
	By \cref{lem:polarization-rigidity}, there is a unique Hermitian inner
	product on \(H_{\mathbb R}\) whose norm-square is the induced quadratic
	form
	\[
		Q_H(\xi)=g(\xi,\xi)
	\]
	and whose imaginary part is determined by the canonical complex
	structure \(J\).
	Thus \(H_{\mathbb R}\), regarded as a complex vector space via \(J\), is
	a pre-Hilbert space.
	Completing \(H_{\mathbb R}\) in the norm induced by this Hermitian form
	produces a complex Hilbert space.
\end{proof}
We denote this completion by \(H\).
% ============================================================
\section{Born rule}
\label{sec:hilbert-born}
% ============================================================
The probabilistic interpretation is now immediate.
\begin{corollary}[Born rule]
	For every \(\xi\in H\),
	\[
		P(\xi)=\|\xi\|^{2}.
	\]
	In particular, the probabilistic scalar channel of the quadratic defect
	sector is exactly the squared norm of the uniquely determined Hermitian
	inner product.
\end{corollary}
\begin{proof}
	By construction of the Hermitian inner product,
	\[
		\|\xi\|^{2}
		=
		\langle \xi,\xi\rangle
		=
		Q_H(\xi).
	\]
	But \(Q_H\) is precisely the quadratic scalar channel induced from the
	unique probability law on first-visible alternatives.
	Therefore
	\[
		P(\xi)=\|\xi\|^{2}.
	\]
\end{proof}
\begin{remark}[Meaning of the norm-square law]
	The Born rule is not postulated here as an external axiom of a Hilbert
	space formalism.
	Rather, the Hilbert norm is reconstructed precisely so that its
	norm-square coincides with the already determined scalar channel on the
	quadratic defect sector.
	Thus the norm-square law is a theorem of the transport structure, not an
	additional probabilistic prescription.
\end{remark}
% ============================================================
\section{Interpretation}
\label{sec:hilbert-interpretation}
% ============================================================
The logic of the chapter may now be summarized without omission.
The closed comparison stack first determines the quadratic transport carrier
\[
	\mathcal K \simeq F^{2}/F^{3}.
\]
The scalar-channel theorem equips that carrier with a positive definite
quadratic law.
Triangle reversal equips it with a unique filtration-compatible
involution.
The ordered orientation double then carries a unique orthogonal complex
structure anti-commuting with total orientation reversal, up to the
harmless sign ambiguity
\[
	J\mapsto -J.
\]
The quadratic scalar law and this quarter-turn determine, by
polarization rigidity, a unique Hermitian inner product.
Completion produces a complex Hilbert space.
Thus the Hilbert framework is not appended externally to the closed
comparison stack.
It is the canonical linear realization of the same intrinsic quadratic
defect layer whose geometric branch, under the inherited
second-jet faithfulness condition, is read on the realized
degree--\(2\) channel attached to the stabilized quadratic carrier,
where nonzero stabilized square classes are equivalent to nonzero
realized curvature.
The chain is
\[
	\begin{aligned}
		 & \text{transport}
		\;\Longrightarrow\;
		F^{2}/F^{3}
		\;\Longrightarrow\;
		\text{scalar law}
		\;\Longrightarrow\;
		\text{orientation double}
		\\
		 & \;\Longrightarrow\;
		J^{2}=-\mathrm{id}
		\;\Longrightarrow\;
		\text{Hermitian rigidity}
		\;\Longrightarrow\;
		\text{Hilbert space}.
	\end{aligned}
\]
The next chapter returns to the parallel geometric branch of the same
intrinsic quadratic defect layer.
Under the inherited second-jet faithfulness condition, that branch is
read on the realized degree--\(2\) channel attached to the stabilized
quadratic carrier, where nonzero stabilized square classes are
equivalent to nonzero realized curvature.
Higher transport jets then begin to contribute genuinely new
geometric information beyond that inherited interface bridge.
\begin{remark}[Parallel realizations of the same carrier]
	The significance of the chapter is not merely that a Hilbert space has
	been reconstructed.
	It is that the Hilbert realization is an intrinsic linear realization
	of the stabilized quadratic defect layer, while the later geometric
	branch is read, under the inherited second-jet faithfulness condition,
	on the realized degree--\(2\) channel attached to the stabilized
	quadratic carrier.
	The subsequent development therefore does not juxtapose two unrelated
	frameworks.
	It develops two linked realizations of one closed comparison
	structure, with the geometric side entering only through that
	inherited interface bridge.
\end{remark}

\section{Conclusion}
\label{sec:hilbert-conclusion}
The chapter proves that Hilbert structure is not auxiliary input: the
Hermitian form, norm-square probability law, and phase sector are all
forced by the stabilized quadratic transport carrier \(F^2/F^3\).
Hilbert realization is therefore an intrinsic determination of the
closed comparison structure, while the later geometric branch is read,
under the inherited second-jet faithfulness condition, on the realized
degree--\(2\) channel attached to the stabilized quadratic carrier.
Accordingly, \cref{ch:einstein} continues from that inherited
interface bridge to the macroscopic compatibility equation carried by
the realized geometric branch.
